\documentclass[12pt]{article}
\usepackage{amsmath, amssymb, amsthm}
\usepackage[margin=1in]{geometry}

\newtheorem{theorem}{Theorem}
\newtheorem{lemma}[theorem]{Lemma}

\title{Problem 273: Sum of Squares}
\author{Project Euler Solution}
\date{}

\begin{document}
\maketitle

\section{Problem Statement}
Let $\mathcal{P} = \{p \text{ prime} : p \equiv 1 \pmod{4},\; p < 150\}$. For each squarefree $N$ that is a product of a non-empty subset of $\mathcal{P}$, let $S(N)$ be the sum of all values $a$ in representations $a^2 + b^2 = N$ with $0 \le a \le b$. Compute $\sum_N S(N)$.

\section{Mathematical Foundation}

\begin{theorem}[Fermat's Two-Square Theorem]
A prime $p$ is representable as $a^2 + b^2$ iff $p = 2$ or $p \equiv 1 \pmod{4}$.
\end{theorem}

\begin{proof}
($\Longleftarrow$) If $p \equiv 1 \pmod{4}$, then $-1$ is a QR mod $p$, so $p \mid (m^2+1)$ for some $m$. In $\mathbb{Z}[i]$, $p \mid (m+i)(m-i)$ but $p \nmid (m \pm i)$, so $p$ is reducible: $p = \pi\bar{\pi}$ with $|\pi|^2 = p$. ($\Longrightarrow$) If $a^2 + b^2 = p$ with $p$ odd, then $(ab^{-1})^2 \equiv -1 \pmod{p}$, so $p \equiv 1\pmod{4}$.
\end{proof}

\begin{theorem}[Gaussian Integer Representation]
Let $N = \prod_{j=1}^k p_j$ be squarefree with each $p_j \equiv 1\pmod{4}$, and $p_j = \pi_j \bar{\pi}_j$ in $\mathbb{Z}[i]$. The representations $a^2+b^2=N$ with $a,b \ge 0$ correspond to
\[
\alpha = u \prod_{j=1}^k \pi_j^{\epsilon_j}\bar{\pi}_j^{1-\epsilon_j}, \quad \epsilon_j \in \{0,1\},\; u \in \{1,i,-1,-i\},
\]
with $(a,b) = (|\operatorname{Re}(\alpha)|,\, |\operatorname{Im}(\alpha)|)$.
\end{theorem}

\begin{proof}
Since $\mathbb{Z}[i]$ is a PID, any $\alpha$ with $|\alpha|^2 = N$ factors uniquely as a unit times a product selecting one from each conjugate pair $\{\pi_j, \bar{\pi}_j\}$. There are $4 \cdot 2^k$ such Gaussian integers; identifying up to sign and coordinate swap yields $2^{k-1}$ distinct representations with $0 \le a \le b$.
\end{proof}

\begin{lemma}[Brahmagupta--Fibonacci Identity]
$(a^2+b^2)(c^2+d^2) = (ac-bd)^2 + (ad+bc)^2 = (ac+bd)^2 + (ad-bc)^2$.
\end{lemma}

\begin{proof}
Direct expansion, or equivalently, the multiplicativity of the Gaussian norm: $|(a+bi)(c+di)|^2 = |a+bi|^2 \cdot |c+di|^2$.
\end{proof}

\begin{theorem}[Incremental Construction]
If $N_1$ has representations $\{(a_i,b_i)\}$ and $p = c^2+d^2$ is a new prime, then the representations of $N_1 p$ are obtained by applying the identity to each $(a_i,b_i)$ with $(c,d)$, yielding two new representations per existing one.
\end{theorem}

\begin{proof}
Follows from the lemma and the Gaussian product $(a_i+b_i\,i)(c \pm d\,i)$.
\end{proof}

\section{Algorithm}
\begin{enumerate}
    \item Identify the 16 primes $p \equiv 1\pmod{4}$ below 150 and their two-square decompositions.
    \item DFS over non-empty subsets, maintaining all $(a,b)$ representations.
    \item At each subset, accumulate $\sum a$ into a global total.
    \item When extending by a new prime, use the Brahmagupta--Fibonacci identity to double the representation list.
\end{enumerate}

\section{Complexity Analysis}
\begin{itemize}
    \item \textbf{Time:} $O(3^{16}) \approx 4.3 \times 10^7$, since a subset of size $k$ contributes $2^{k-1}$ representations and $\sum_{k=1}^{16}\binom{16}{k}2^{k-1} = (3^{16}-1)/2$.
    \item \textbf{Space:} $O(2^{16})$ for the largest representation list at DFS depth 16.
\end{itemize}

\section{Answer}

\[
\boxed{2032447591196869022}
\]

\end{document}
